\documentclass[11pt]{letter}
\usepackage[margin=1in]{geometry}
\usepackage{hyperref}

\begin{document}

\begin{letter}{
Editor-in-Chief \\
Neurocomputing \\
}

\opening{Dear Editor,}

We submit our manuscript \textbf{``Higher Rank, More Hallucination? How LoRA Capacity Controls Label Compliance in Fine-Tuned LLMs''} for consideration in \emph{Neurocomputing}.

\textbf{Summary.} We investigate how LoRA rank---the primary capacity control in parameter-efficient fine-tuning---affects label compliance. Testing ranks 16, 32, 64, and 128 across three models, we discover a non-monotonic relationship: moderate ranks (32--64) maximize compliance, while both lower and higher ranks degrade it. Higher LoRA rank activates more pre-training knowledge, which paradoxically introduces competing label vocabularies that cause hallucination.

\textbf{Key contributions:}
\begin{enumerate}
\item A rank--compliance curve showing non-monotonic behavior: rank 64 outperforms both rank 16 (underfitting) and rank 128 (over-activation).
\item A formal definition of Rank Sensitivity that quantifies how a model's compliance changes with capacity.
\item Practical LoRA rank selection guidelines for schema-constrained classification tasks.
\end{enumerate}

\textbf{Relevance to Neurocomputing.} LoRA is the dominant method for adapting large neural networks to domain tasks. Our finding that rank selection has non-obvious effects on output compliance provides essential guidance for the Neurocomputing community working on parameter-efficient adaptation.

\textbf{Novelty.} Prior LoRA studies focus on accuracy and convergence speed. We are the first to show that LoRA rank has a direct, non-monotonic effect on label hallucination, and that this effect is model-dependent.

This manuscript has not been published previously and is not under consideration elsewhere.

\closing{Respectfully submitted,\\[2em]
Nutthakorn Chalaemwongwan\\
Department of Computer Engineering, KOSEN-KMITL\\
King Mongkut's Institute of Technology Ladkrabang\\
nutthakorn.ch@kmitl.ac.th
}

\end{letter}
\end{document}
